\section{The Step to Coupled PDEs}\label{sec:coupled}

We now couple species to allow for multi-linear systems, a direct expansion from the canon work we expand on!

\subsection{Reaction-Diffusion Family and OTS}

For $N$ species with diagonal diffusion, linear reaction coupling, and (here) zero noise --- the deterministic skeleton of the moment systems to come ---
\begin{equation}\label{eq:coupledrd}
\partial_t u^{(i)} = D_i\,\partial_{xx}u^{(i)} + \sum_j R_{ij}u^{(j)}.
\end{equation}

\begin{lemma}[OTS-NIDC for fixed indices]
  For equation \eqref{eq:coupledrd}, the optimal step for each $u^{(i)}$ individually, as derived in the single variable case, is exactly $\dto^{(i)} = \frac{\dx^2}{6D_i}$, with first-level corrections $C^{(1),(i)}$ of the same algebraic form as \eqref{eq:C1}, derived term-by-term in \texttt{OTS-NIDC.md} \S8--9.
\end{lemma}

With this, we get the desired theorem as follows.

\begin{theorem}[First-level accuracy for fixed indices]\label{thm:cpl-p1}
  The scheme

  \begin{equation}\label{eq:coupledscheme}
      u^{(i), n+1}_{j} = u^{(i), n}_{j} + \dto^{(i)} \left[ D_i \ddx u^{(i), n}_{j} + \sum_{\ell} R_{i\ell} u^{(\ell), n}_{j} \right] + C^{(1), (i)}
  \end{equation}

  has similar step error and global order as given by Theorem \ref{thm:p1}.
\end{theorem}

\subsubsection{Competing Time Steps and Subcycling}

Notice that the scheme presented by Theorem \ref{thm:cpl-p1} relies on having $u^{(\ell), n}_{j}$, the value of the separate PDE $u^{(\ell)}$ at time-step $n$. However, since each equation carries its own optimal step, the required value is not readily available whenever the diffusivities differ, $D_i \neq D_j$ for some $i \neq j$, i.e.\ whenever no single $\dto$ serves the whole system.

The two failure modes of a naive single step frame the problem.

\paragraph{Failure 1: shared step without correction.} If species 1 (fast,
$D_1$, optimal step $\tau_1=\dto^{(1)}$) is forced onto species 2's optimum
$h=\tau_2=\dto^{(2)}$, its Group-1
cancellation breaks. The one-step identity
\begin{equation}\label{eq:mismatchid}
e(h) \;=\; h D_1\,\tilde u_{xxxx}
\Bigl[\tfrac{h D_1}{2} - \tfrac{\dx^2}{12}\Bigr] + O(h^3)
\end{equation}
is exact for every $h$ --- where $\tilde u_{xxxx}=\partial_x^4u^{(1)}$ is evaluated on the exact solution at the step, the representative of the LTE remainder as in \S\ref{sec:method} --- and writing $h = \tau_1 + \varepsilon$ turns
the bracket into $\varepsilon D_1/2$ \emph{identically} (no Taylor truncation in
$\varepsilon$). 

We remain, however, that this stencil is not ideal for each step, and it leads to questions of stability. In fact,

\paragraph{Failure 2: naive splitting.} Lie--Trotter splitting with the
coupling frozen over a macro step $H$ incurs a sustained $O(H)$ forcing error;
accumulated over $T/H$ cycles this is a global $O(H) = O(\dx^2)$ error ---
worse than the uncorrected base scheme. Freezing is not viable; coupling
data must be interpolated. 

The rest of this paper introduces methods of attempting to corroborate the differences between these time-steps in ways to recover the accuracy that we would lose by just picking one or the other. 
